% !TEX root = ../main.tex
\section{Benchmark Suite}
\label{sec:evaluation}

\textsc{anneal} ships with a benchmark suite of five optimization targets spanning three domains (prompt optimization, code optimization, and documentation generation) that exercise both evaluation modes (deterministic and stochastic) and a range of optimization objectives.
The suite is part of the open-source distribution and is intended as a reference workload for end-to-end configuration comparison.
The mechanism-level validation in \cref{sec:gate-false-positives,sec:gate-sa,sec:gate-retrieval} uses synthetic benchmarks with analytically known ground truth to isolate individual subsystems, and the end-to-end comparison in \cref{sec:end-to-end} runs the full benchmark suite under the four-configuration protocol described below.
This section documents the suite's targets and the statistical protocol applied to the end-to-end run.

\subsection{Benchmark Targets}
\label{sec:benchmark-suite}

\cref{tab:benchmark-targets} summarizes the five benchmark targets.
Each target defines an artifact, an evaluation function, and a set of editable files; the evaluation function and its criteria are immutable and outside the agent's mutation scope.

\begin{table}[t]
  \centering
  \caption{Benchmark targets. \emph{Eval}: D = deterministic (single scalar), S = stochastic ($N$ samples $\times$ $K$ binary criteria). \emph{Dir}: optimization direction.}
  \label{tab:benchmark-targets}
  \small
  \begin{tabular}{@{}lllcrl@{}}
    \toprule
    ID & Target & Domain & Eval & $K$ & Dir \\
    \midrule
    B1 & Classification prompt & Prompt & S & 5 & max \\
    B2 & Summarization prompt & Prompt & S & 4 & max \\
    B3 & Performance opt. & Code & D & --- & min \\
    B4 & Code quality refactor & Code & D & --- & max \\
    B5 & README generation & Docs & S & 6 & max \\
    \bottomrule
  \end{tabular}
\end{table}

The targets were selected to satisfy three coverage criteria.
First, they span both evaluation modes: B3 and B4 use deterministic evaluation (a single scalar from a shell command: wall-clock time and a composite quality score, respectively), while B1, B2, and B5 use the $N \times K$ stochastic evaluation framework with binary criteria.
Second, they cover three distinct optimization domains (prompt engineering, code transformation, and technical writing) so that results are not tied to a single domain.
Third, they exercise different optimization directions: maximization (B1, B2, B4, B5) and minimization (B3), testing the engine's directional generality.

Each target ships with a TOML configuration specifying artifact paths, scope constraints, evaluation commands, and scoring criteria.
The evaluation functions are self-contained: B3 runs a subprocess and parses wall-clock time; B4 runs a composite quality scorer (covering correctness, input validation, error handling, and security); B1, B2, and B5 invoke an LLM judge with fixed rubric criteria and aggregate binary pass/fail scores across $N = 5$ test prompts.

The harness ships four named configurations (raw, greedy, control, treatment) defined in \texttt{benchmarks/suite/runner.py} for end-to-end comparison.
\Cref{tab:e2e-configs} documents the intended configuration semantics, and \cref{sec:end-to-end} reports the empirical comparison.

\subsection{Result Provenance}
\label{sec:result-provenance}

The end-to-end tables in \cref{sec:end-to-end} are regenerated from the archived JSONL records in \texttt{benchmarks/raw\_results\_seeds1-5/} and the generated analysis outputs in \texttt{benchmarks/results\_seeds1-5/}.
The archived records are the authoritative empirical source for the reported end-to-end numbers.
The current benchmark runner and analysis scripts are included with the artifact, but the archived run was produced under mixed historical model routing; therefore the end-to-end results should be read as a result bundle with explicit provenance, not as a claim that a fresh checkout will reproduce identical LLM trajectories.

\begin{table*}[t]
  \centering
  \caption{Provenance for the archived end-to-end result bundle used in \cref{sec:end-to-end}.}
  \label{tab:result-provenance}
  \footnotesize
  \begin{tabular}{@{}p{3.0cm}p{13.0cm}@{}}
    \toprule
    Field & Value \\
    \midrule
    Analysis commit & \texttt{0b48517de4cff5f87c4697a149b78dc2c955bfe0} \\
    Raw records & \texttt{benchmarks/raw\_results\_seeds1-5/} \\
    Analysis outputs & \texttt{benchmarks/results\_seeds1-5/} \\
    Seeds & $1,2,3,4,5$ \\
    Observed mutation models & \texttt{claude-sonnet-4-6}, \texttt{gpt-5.4}, \texttt{gemini-3.1-pro-preview} \\
    Drafts per record & $1$ \\
    Analysis command & \texttt{uv run python benchmarks/analysis/run\_analysis.py --results-dir benchmarks/raw\_results\_seeds1-5 --output-dir benchmarks/results\_seeds1-5} \\
    \bottomrule
  \end{tabular}
\end{table*}

This provenance constraint does not affect the mechanism gates, which are deterministic local benchmarks.
It does constrain interpretation of the end-to-end study: the archived records support descriptive, paired-seed comparisons and cost telemetry under the recorded execution environment, but they do not establish exact replayability of historical LLM outputs from the current repository state.
The current runner defines intended enhancement flags for the treatment configuration, but the archived result records are the empirical source of truth; consequently, this paper does not claim feature-level attribution from the treatment bundle.

\subsection{Statistical Design for End-to-End Comparison}
\label{sec:stat-design}

The harness implements a paired-seed experimental design for comparing configurations.
Each configuration runs 5 independent trials per benchmark target, each with a distinct random seed that controls LLM sampling temperature jitter and stochastic evaluation order.
Every trial was budgeted for 50 experiments (mutations), yielding a maximum planned budget of $5 \text{ targets} \times 4 \text{ configs} \times 5 \text{ runs} \times 50 \text{ experiments} = 5{,}000$ experiments for the end-to-end comparison reported in \cref{sec:end-to-end}; deterministic targets can stop early when a perfect score is reached.
The harness supports larger seed counts; we report the archived $n = 5$ run because no additional benchmark experimentation was performed.
At this $n$, the Wilcoxon signed-rank test has a two-sided $p$-value floor of $0.0625$, which prevents conventional per-target significance for perfect five-seed directional sweeps and makes the corrected test deliberately conservative (\cref{sec:end-to-end-significance}).

\paragraph{Pairwise comparison.}
For each pair of configurations on each target, the harness applies the Wilcoxon signed-rank test across the 5 paired final scores (paired by random seed).
This nonparametric test is appropriate because score distributions are not assumed normal, and pairing by seed controls for LLM sampling variability.

\paragraph{Multiple comparison correction.}
Holm--Bonferroni correction is applied across the five targets for each configuration pair, controlling the family-wise error rate at $\alpha = 0.05$.
This mirrors the correction used within the engine itself for consolidation-window decisions (\cref{sec:stat-comparison}).

\paragraph{Effect size.}
The harness reports Cohen's $d$ with 95\% bootstrap confidence intervals (10{,}000 resamples, bias-corrected and accelerated) using the bootstrap methodology of Efron and Tibshirani~\cite{efron1994bootstrap}.
Bootstrap resampling uses deterministic seeding with float-precision normalization for cross-platform reproducibility.

\paragraph{Cost telemetry.}
The analysis reports total recorded USD cost per run using the cost fields stored in the raw JSONL records.
These costs are telemetry for the archived model-routing setup in \cref{tab:result-provenance}; they should not be interpreted as model-independent cost laws.
